% !TEX root = ../prise-en-main-medical-vrp.tex

\section{Votre journée type}
\label{sec:quotidien}

\subsection{Planifier une intervention}

\begin{workflowstep}[Agenda]
  \menu{Agenda} → vues \textbf{mois} ou \textbf{semaine}. Sélectionnez \structure{}, \seance{} et \patient{}. En mode \textbf{Édition}, cliquez un jour (08{,}00) ou un créneau horaire, ou ouvrez le formulaire. Dupliquez une séance vers demain, la semaine suivante ou une date choisie.
\end{workflowstep}

\screenshot[keepaspectratio]{05-agenda.png}{Agenda — sélection structure, séance, patient}

Chaque créneau planifié alimente le montant \textbf{non facturé TTC} visible sur la liste des \structure{}s.

\subsection{Préparer la facturation}

La facturation se fait \textbf{sur la fiche \structure{}}, pas dans l'agenda.

\begin{workflowstep}[Facturation mensuelle]
  Accueil → cliquez sur la \structure{} → section \textbf{Facturation} (\texttt{\#billing}). Naviguez mois par mois. Vérifiez les \seance{}s, \patient{}s et montants HT/TTC de chaque ligne. Créez ou rattachez une facture (mode Édition).
\end{workflowstep}

\screenshot[keepaspectratio]{06-facturation-structure.png}{Préparation facturation sur la fiche structure}

Bouton utile~: \menu{Sauter aux sessions non facturées} pour retrouver le dernier mois en attente.

\subsection{Suivre les encaissements}

\begin{workflowstep}[Trésorerie]
  \menu{Trésorerie} → synthèse mensuelle (émis, payé, planifié). Importez un relevé \texttt{.xlsx} (onglet \menu{Banque}) et rapprochez les crédits avec vos factures.
\end{workflowstep}

\screenshot[keepaspectratio]{07-tresorerie.png}{Synthèse trésorerie}

Un rapprochement bancaire renseigne automatiquement la date de paiement de la facture si elle était vide.

\subsection{Documents}

Joignez conventions, ordonnances scannées ou bons de commande sur la fiche \structure{} (\menu{Documents}) pour centraliser le dossier administratif.
